\section{Failure: a contracting mechanism and local realizations}
\label{art:failure}

The positive theorems identify geometry that permits convergence. We now show
why convexity and regularity of the branch formulas cannot supply an
unrestricted deterministic replacement. The first tempting assertion is
global: compact sublevels might seem enough to prevent convergence to a wrong
corner. The second is local: for one finite maximum of very smooth convex
branches, perhaps every sufficiently close start succeeds. We construct an
actual failure of each assertion, and then examine what changes when positive
lower curvature is imposed.

The mechanism has three parts. A three-step BFGS cycle returns the inverse
matrix to a smaller multiple of itself while approaching a nonstationary
corner. Strict function and gradient inequalities force the specified line
search to produce that cycle. Finally, local copies of the corner accumulate
at the true minimizer of one objective. Matching spatial and objective scales
keeps the initial matrix fixed across the copies. The distinction between a
formal cycle and an implemented run matters at every part of this argument.

\subsection{The specified implementation and the local guarantee at issue}
At a differentiable point $x$, let $g=\nabla F(x)$ and $p=-Hg$.
The line search starts with lower bracket $\alpha_{\rm lo}=0$, upper bracket
$\alpha_{\rm hi}=+\infty$, and trial length $\alpha=1$. It tests Armijo first and weak
Wolfe second:
\begin{equation}\label{art:failure_tests}
 F(x+\alpha p)\le F(x)+c_1\alpha g^Tp,
 \qquad \nabla F(x+\alpha p)^Tp\ge c_2g^Tp.
\end{equation}
An Armijo rejection sets $\alpha_{\rm hi}=\alpha$; a Wolfe rejection sets $\alpha_{\rm lo}=\alpha$.
The next length is $(\alpha_{\rm lo}+\alpha_{\rm hi})/2$ when $\alpha_{\rm hi}<\infty$, and $2\alpha_{\rm lo}$ otherwise.
If both tests hold, the length is accepted. This is the unit-first
bracketing implementation of Lewis--Overton~\cite[Algorithm~4.6]{LO}.
All comparisons in the construction are strict, so strict acceptance tests
give the same decisions.

Write the full inverse update as
\begin{equation}\label{art:failure_update}
 \mathcal B(H,g,g^+,\alpha)
 =\left(I-\frac{sy^T}{y^Ts}\right)H
  \left(I-\frac{ys^T}{y^Ts}\right)+\frac{ss^T}{y^Ts},
 \quad s=-\alpha Hg,\quad y=g^+-g.
\end{equation}
No resetting or rescaling of the stored matrix is performed.

For this routine, even identity initialization does not give a neighborhood
in which every differentiable start succeeds under the following regularity
assumptions.
\begin{theorem}[Local failure with globally smooth convex branches]
\label{art:local-failure}
There are four globally convex functions $f_i\in C^{2,1}(\R^2)$ such that
\[
 0\preceq\nabla^2f_i(x)\preceq64I\qquad(i=0,1,2,3; x\in\R^2)
\]
and $F=\max_i f_i$ has compact sublevel sets and a unique nonsmooth
minimizer at zero. For this one objective there are fixed starting points
$x_0^{(\ell)}\to0$, all with $H_0=I$, such that, for every fixed pair
\[
 0<c_1\le\tfrac13,\qquad c_1<c_2<1,
\]
the deterministic unit-first bracketing implementation of Section~\ref{art:failure} produces infinite full-BFGS runs satisfying
\[
 x_k^{(\ell)}\to \bar x_\ell\ne0,\qquad
 F(x_k^{(\ell)})\downarrow F(\bar x_\ell)>F(0)=0,\qquad
 0\notin\partial F(\bar x_\ell).
\]
Every queried point is differentiable. Each line search takes at most five
trials, and its accepted lengths repeat $1,1/4,1/16$.
The objective, initial points, initial matrix and complete decision tree
are the same throughout the displayed parameter range. The supremum of
the norms of all accepted and rejected query points in the $\ell$th run
tends to zero as $\ell\to\infty$.
\end{theorem}


The objective and all starts are chosen before the displayed line-search
parameters. The failure set supplied by the theorem is a sequence of starts;
it is not a positive-measure basin. There is no uniform positive lower
curvature requirement on these branches.

\subsection{Designing a cycle around a wrong corner}
We seek a corner with a strict descent direction that the cycle nevertheless
fails to follow to the minimum. Three slopes in the plane provide the needed
freedom. Begin with $H_0=I$ and fix the first two slopes so that the unit step
has a simple exact update. The third slope has two free coordinates. We choose
them to close the matrix cycle at a scalar multiple $\rho I$ and choose the
starting point to close the displacement cycle at $\rho z_0$. For the maximum
of linear functions, those two identities allow the same geometry to repeat
at each smaller scale.

The target lengths $1,1/4,1/16$ fit the bisections of the unit-first search.
Choosing them in the algebra does not prove that the algorithm accepts them:
we must also exclude every earlier candidate and verify both tests at the
intended endpoint. The strict trial certificate below supplies exactly this
missing dynamical condition.

Let $(c_\circ,d_\circ)$ be the algebraic pair specified in
Appendix~\ref{art:seed_certificate}. Define
\[
 g_0=(1,0)^T,\qquad g_1=(-1,-1/8)^T,\qquad g_2=(c_\circ,d_\circ)^T,
 \quad G(z)=\max_{0\le i\le2}g_i^Tz.
\]
The useful coarse bounds are $0.318<c_\circ<0.319$ and
$-1.026<d_\circ<-1.025$. Put
\begin{equation}\label{art:seed_rho}
 \rho=\frac{129(c_\circ-8d_\circ)}{(16c_\circ-127d_\circ-16)(127c_\circ+16d_\circ+129)},
 \qquad 0.060<\rho<0.061.
\end{equation}
Set
\begin{align}
 M_0&=I,\quad
 M_1=\begin{pmatrix}129/256&-1/16\\-1/16&1\end{pmatrix},
 &M_2&=\mathcal B(M_1,g_1,g_2,1/4),\label{art:seed_matrices}\\
 s_0&=(-1,0)^T,\quad s_1=(127/1024,1/64)^T,
 &s_2&=\rho(1-c_\circ,-d_\circ)^T.\label{art:seed_steps}
\end{align}
The point is then determined by the displacement closure
$s_0+s_1+s_2=-(1-\rho)z_0$; explicitly, let
\begin{equation}\label{art:seed_points}
 z_0=\frac1{1-\rho}
 \begin{pmatrix}897/1024-\rho(1-c_\circ)\\ \rho d_\circ-1/64\end{pmatrix},
 \qquad z_{j+1}=z_j+s_j\quad(j=0,1,2).
\end{equation}

\begin{lemma}[Exact cycle and strict trial decisions]\label{art:seed_cycle}
The matrices in \eqref{art:seed_matrices} are positive definite, and
\[
 \mathcal B(I,g_0,g_1,1)=M_1,\qquad
 \mathcal B(M_2,g_2,g_0,1/16)=\rho I,\qquad z_3=\rho z_0.
\]
At $z_j$, $j=0,1,2$, the unique active gradient is $g_j$.
With $p_j=-M_jg_j$, the entire line-search query list is
Table~\ref{art:trial_table}. At every rejection the objective strictly
exceeds its current value. At every acceptance,
\begin{equation}\label{art:seed_margins}
 G(z_j)+\tfrac13g_j^Ts_j-G(z_{j+1})>0,\qquad
 g_{j+1}^Ts_j>0>g_j^Ts_j,
\end{equation}
where $g_3=g_0$. All listed queries and the three starting points have a
unique active branch and satisfy
\begin{equation}\label{art:seed_query_bounds}
 |g_i^Tz|<3 (i=0,1,2),\qquad [z]_2<2,\qquad G(z)>0.
\end{equation}
\end{lemma}

\begin{table}[t]
\centering
\setlength{\tabcolsep}{5.8pt}
\begin{tabular}{cccccl}
\toprule
Phase $j$ & \shortstack{$10^4G$\\current} & Trial $\alpha$
 & \shortstack{$10^4G$\\trial} & Active branch & Decision\\
\midrule
0 & $(8884,8885)$ & $1$ & $(1218,1219)$ & 1 & accept\\
\midrule
1 & $(1218,1219)$ & $1$ & $(3845,3846)$ & 0 & reject\\
  &               & $1/2$ & $(1364,1365)$ & 0 & reject\\
  &               & $1/4$ & $(723,724)$ & 2 & accept\\
\midrule
2 & $(723,724)$ & $1$ & $(6687,6691)$ & 0 & reject\\
  &             & $1/2$ & $(3406,3408)$ & 0 & reject\\
  &             & $1/4$ & $(1765,1767)$ & 0 & reject\\
  &             & $1/8$ & $(944,946)$ & 0 & reject\\
  &             & $1/16$ & $(534,536)$ & 0 & accept\\
\bottomrule
\end{tabular}
\caption{The deterministic trial tree, read against each phase's current
value. ``Current'' means $G(z_j)$ and ``trial'' means
$G(z_j+\alpha p_j)$. Every rejected trial raises the objective and
therefore fails Armijo. The accepted rows also satisfy the strict
Armijo and weak-Wolfe margins in \eqref{art:seed_margins}.
All displayed intervals are rational enclosures, not rounded values.}
\label{art:trial_table}
\end{table}

The proof of the lemma, including the root isolation, matrix closure and
rational certificate, is in Appendix~\ref{art:seed_certificate}. The logic
of the search can already be read from its margins. Every rejected candidate
has a value larger than the current value, so it fails Armijo for every
$c_1>0$. At each accepted candidate the decrease exceeds one third of the
predicted decrease, while the new directional derivative is positive and
the old one is negative. Hence Armijo holds for $c_1\le1/3$ and weak Wolfe
holds for every $c_2>0$. Thus the displayed lengths are outputs of the
specified search for the entire parameter range in
Theorem~\ref{art:local-failure}; no accepted-length sequence is assumed.
The routine never needs a Wolfe rejection on this example.

Figure~\ref{art:cycle-figure} displays this unit-first geometry: phase zero
accepts its first query, while the next two phases bisect past the rejected
candidates before reaching their accepted points. Panel (b) includes all nine
queries; panel (c) separates the accepted points near the wrong corner. Line styles distinguish the three phases
in grayscale, and the accepted lengths are labeled explicitly.

\begin{figure}[t]
\centering
\includegraphics[width=\textwidth]{figures/cycle_trials.pdf}
\caption{The accepted three-step cycle, every queried candidate, and an
enlargement at the nonstationary origin. Solid, dashed and dash--dot rays
mark phases $0,1,2$, whose accepted lengths are $1,1/4,1/16$.
The panels use different horizontal and vertical scales, as indicated.
Coordinates are rendered from the rational center of the certified root
box; exact cycle closure and trial decisions follow from
Appendix~\ref{art:seed_certificate}. This drawing is not evidence of
stable-manifold membership for the perturbed construction.}
\label{art:cycle-figure}
\end{figure}

The finite calculation becomes an infinite run through homogeneity. A smaller
copy has the same active slopes, and its inverse matrix produces proportionally
smaller steps. Both sides of every line-search comparison have the same scale.
The two cycle closures therefore give the exact induction
\begin{equation}\label{art:cone_orbit}
 x_{3m+j}=\rho^mz_j,\qquad H_{3m+j}=\rho^mM_j,\qquad j=0,1,2.
\end{equation}
Each query in phase $j$ is $\rho^m$ times its base query. Gradients are
unchanged, function differences and gradient--step products scale by
$\rho^m$, and \eqref{art:failure_update} scales the updated matrix by
$\rho^m$. Thus every branch choice and every rejection is preserved, as
well as acceptance. The iterates converge to zero, but the direction
$w=(-1,16)^T$ satisfies
\[
 g_0^Tw=g_1^Tw=-1,\qquad g_2^Tw=-c_\circ+16d_\circ<-1.
\]
Hence $G(tw)=-t$ for $t\ge0$, so this limit is nonstationary.

\begin{corollary}[Coercive polyhedral global failure]\label{art:failure_polyhedral}
The same deterministic run, with $H_0=I$, fails for
\[
 \overline G(u,v)=\max\{u,-u-v/8,c_\circ u+d_\circ v,v-17\}.
\]
This function is coercive and has unique minimizer $(-1,16)$, with value
$-1$, while the run converges to the nonstationary origin with value zero.
\end{corollary}
\begin{proof}
At every query, \eqref{art:seed_query_bounds} and $0<\rho<1$ give
$v-17<0<G(u,v)$, so the fourth branch is inactive. The convex combination
of branches $0,1,3$ with weights $8/17,8/17,1/17$ is identically $-1$.
Equality in this lower bound forces all three branches to equal $-1$,
hence $(u,v)=(-1,16)$; the remaining branch is strictly smaller there.
If $\overline G(u,v)\le t$, then
$u\le t$, $v\le t+17$ and $-u-v/8\le t$, which also give lower bounds
on both coordinates. Every sublevel is therefore compact.
\end{proof}
This already excludes global convergence from every start for coercive finite
polyhedral objectives under this routine. Its wrong corner is away from the
optimizer, so it is consistent with the local finite-partition convergence
theorem. To rule out a local neighborhood guarantee, we need bad corners
arbitrarily near the optimizer of the \emph{same} function.

\subsection{Matching the spatial scale, objective scale and initial matrix}
A change of scale also changes which inverse matrix represents the same
search direction. This simple covariance explains both local constructions.
Suppose, near the current and queried points under consideration, that
\begin{equation}\label{art:matched-local-model}
 F(\bar x+\lambda z)=C+\sigma\mathcal G(z),\qquad \lambda,\sigma>0.
\end{equation}
Differentiation gives $g_F=(\sigma/\lambda)g_{\mathcal G}$. Consequently,
for the same candidate length, the corresponding inverse matrix and step are
\begin{equation}\label{art:matched-scaling}
 H_F=\frac{\lambda^2}{\sigma}H_{\mathcal G},\qquad
 p_F=\lambda p_{\mathcal G},\qquad s_F=\lambda s_{\mathcal G}.
\end{equation}
At an accepted pair $y_F=(\sigma/\lambda)y_{\mathcal G}$. Thus
$y_F^Ts_F=\sigma y_{\mathcal G}^Ts_{\mathcal G}$; the two projection
factors in the inverse update are identical, and its rank-one term is
multiplied by $\lambda^2/\sigma$. This proves the matrix relation in
\eqref{art:matched-scaling} again after the update. Function differences
and all gradient--step products are multiplied by $\sigma>0$, so the
Armijo and weak-Wolfe comparisons, including their strict signs, agree.

This is a covariance calculation, not an existence theorem for a local
trajectory. To transfer a particular run, the local model must cover the
current point and every candidate the search actually visits. Under that
containment, induction on the bracket decisions transfers the full trial
tree, and induction on accepted steps transfers the matrices. We verify
that containment separately in each construction below.

The choice $\sigma=\lambda^2$ is decisive: it leaves the initial matrix
unchanged. In the first embedding, the corner's objective height is $a_\ell^3$
and its affine slopes are $a_\ell^2g_i$. Taking spatial scale
$\lambda=a_\ell^2$ gives $\sigma=a_\ell^4=\lambda^2$, so every copy starts
with $I$. Its affine neighborhood has width proportional to $a_\ell$, much
larger than the entire copied trial tree, whose width is proportional to
$a_\ell^2$. This separation of scales makes all-query containment possible.

\subsection{Convex interpolation that leaves room for every query}
We now construct the scalar functions that implement this design. Their node
values are $a_\ell^3$ and their node slopes are $a_\ell^2$. A monotone derivative
connects neighboring nodes and is constant on a neighborhood of each node.
The constant regions create the local affine corner; the cubic transition
makes the second derivative continuous and globally Lipschitz even as the
nodes accumulate at zero. The following formulas fix the construction
completely.

Let $\theta=1/16$, $a_\ell=\theta^\ell$, $\varepsilon=1/128$, and
$b_*=2c_\circ-32d_\circ\in(33,34)$. For $b\in\{2,b_*\}$ define
\[
 t_b=\frac{1-(1+\theta+\theta^2)/b}{1-\theta^2},\qquad
 v=\frac{\xi-t_b+\varepsilon}{2\varepsilon}.
\]
\[
 P_b(\xi)=
 \begin{cases}
 0&\xi\le t_b-\varepsilon,\\
 3v^2-2v^3&|\xi-t_b|<\varepsilon,\\
 1&\xi\ge t_b+\varepsilon,
 \end{cases}
 \quad J_b(\xi)=\int_0^\xi P_b(t)\,dt.
\]
The middle formula for $J_b$ is $2\varepsilon(v^3-v^4/2)$; its outer
formulas are zero and $\xi-t_b$. For
$b\theta a_\ell\le x\le ba_\ell$, set
\begin{equation}\label{art:upper_scalar}
 \varphi_b(x)=(\theta a_\ell)^3+
 ba_\ell^3(1-\theta)
 \{\theta^2\xi+(1-\theta^2)J_b(\xi)\},\qquad
 \xi=\frac{x/(ba_\ell)-\theta}{1-\theta}.
\end{equation}
Complete the definition by
\[
 \varphi_b(x)=
 \begin{cases}
 1-3(x+1)+3(x+1)^2&x<-1,\\
 -x^3&-1\le x\le0,\\
 1+x-b&x\ge b.
 \end{cases}
\]

\begin{lemma}[Scalar interpolation]\label{art:upper_scalar_lemma}
The functions $\varphi_b$ are convex, coercive and globally $C^{2,1}$,
are nonnegative, and vanish only at zero. They satisfy
\begin{align}
 0&\le\varphi_b''\le\max\{6,102/b\},\label{art:scalar_curvature}\\
 \varphi_b(ba_\ell+h)&=a_\ell^3+a_\ell^2h
       \qquad(|h|\le a_\ell/2).\label{art:scalar_plateau}
\end{align}
Their second derivatives are globally Lipschitz with constant
\[
 L_b=\max\{6,139264/(5b^2)\}.
\]
\end{lemma}
\begin{proof}
The ramp is monotone, $C^1$, and satisfies
$0\le P_b'\le96$ and $\int_0^1P_b=1-t_b$. On a positive interval,
\[
 \varphi_b'=a_\ell^2\{\theta^2+(1-\theta^2)P_b(\xi)\},\qquad
 \varphi_b''=\frac{a_\ell(1+\theta)}bP_b'(\xi).
\]
The identity
$\theta^2+(1-\theta^2)(1-t_b)=(1+\theta+\theta^2)/b$
gives matching values $\varphi_b(ba_\ell)=a_\ell^3$. The derivative is
constant near both ends of each interval, so the first two derivatives
also match, with values $a_\ell^2$ and zero at $ba_\ell$.

As $x\downarrow0$, the bounds $a_\ell\le x/(b\theta)$ on its interval
give $\varphi_b(x)=O(x^3)$, $\varphi_b'(x)=O(x^2)$ and
$\varphi_b''(x)=O(x)$. These match the negative cubic at zero.
The formulas also match through second order at $-1$ and $b$.
The derivative is negative on the negative half-line and positive on the
positive half-line; the outer formulas diverge to infinity. This proves
convexity, coercivity and the minimum assertions.

The constant-derivative lengths below and above $ba_\ell$, divided by $a_\ell$,
are
\[
 b(1-\theta)(1-t_b-\varepsilon),\qquad
 \frac{b(1-\theta)(t_b-\varepsilon)}\theta.
\]
Both exceed $1/2$ for $2\le b\le34$, by substitution in the displayed
rational formulas; the affine outer extension covers the upper side when
$\ell=0$. This proves \eqref{art:scalar_plateau}.
Finally $|P_b''|\le3/(2\varepsilon^2)$ in each ramp, so differentiation
once more gives $|\varphi_b'''|\le139264/(5b^2)$ there. The negative
cubic has third derivative of magnitude six. Continuity of
$\varphi_b''$ permits summing the Lipschitz estimates over finitely many
joins away from zero; taking limits at zero covers the infinitely many
accumulating joins and gives the global bound.
\end{proof}

\subsection{One objective and a uniform induction over all queries}
Define
\begin{equation}\label{art:upper_objective}
 F(u,v)=\max\{\varphi_2(u),\varphi_2(-u-v/8),
                    \varphi_{b_*}(c_\circ u+d_\circ v),v\}.
\end{equation}
The Hessian of each of the first three branches is a nonnegative scalar
multiple of $g_ig_i^T$. Equation~\eqref{art:scalar_curvature} gives the
upper bounds $51$, $51(65/64)<52$, and $12$, respectively, on their
operator norms; the last branch has zero Hessian. Their Hessians are
globally Lipschitz because $\varphi_b''$ is Lipschitz and the arguments
are linear. The first two scalar arguments are independent, so their
coercivity makes every sublevel of $F$ compact. Nonnegativity and their
unique scalar zeros show $F(x)=0$ exactly at $x=0$. At zero the first
three gradients vanish and the last is $(0,1)^T$, giving
$F'(0;h)=\max\{0,h_2\}$; thus the minimizer is nonsmooth.

Let $\bar x_\ell=a_\ell(2,-32)^T$. The first two scalar arguments at $\bar x_\ell$ equal
$2a_\ell$, and the third equals $b_*a_\ell$. Consequently
\begin{equation}\label{art:upper_local_model}
 F(\bar x_\ell+h)=a_\ell^3+a_\ell^2G(h)
 \quad\hbox{if } |g_i^Th|\le a_\ell/2 (i=0,1,2).
\end{equation}
Indeed, the first three branches are on their affine regions and each
has value at least $a_\ell^3/2$. The first two slab inequalities imply
$|h_1|\le a_\ell/2$ and $|h_2|\le8a_\ell$, so the fourth branch is negative.
In particular $F(\bar x_\ell)=a_\ell^3$ and
$F'(\bar x_\ell;(-1,16)^T)=-a_\ell^2<0$.

\begin{proof}[Theorem~\ref{art:local-failure}]
For $\ell\ge1$ take $x_0^{(\ell)}=\bar x_\ell+a_\ell^2z_0$ and $H_0=I$. We claim that
\begin{equation}\label{art:upper_orbit}
 x_{3m+j}^{(\ell)}=\bar x_\ell+a_\ell^2\rho^mz_j,\qquad
 H_{3m+j}^{(\ell)}=\rho^mM_j,\qquad j=0,1,2.
\end{equation}
At a current point of this form its gradient is $a_\ell^2g_j$. Every trial
on the finite candidate list therefore has the form
$\bar x_\ell+a_\ell^2\rho^mz_{\rm trial}$. Equation~\eqref{art:seed_query_bounds}
gives
\[
 |g_i^T(a_\ell^2\rho^mz_{\rm trial})|<3a_\ell^2
 \le3a_\ell/16<a_\ell/2.
\]
Thus \eqref{art:upper_local_model} is valid at every candidate query,
not only the final accepted one. Unique activity is preserved, and
\[
 F(\bar x_\ell+a_\ell^2\rho^mz_{\rm trial})
   =a_\ell^3+a_\ell^4\rho^mG(z_{\rm trial}),\qquad
 \nabla F=a_\ell^2g_{\rm active}.
\]
These identities instantiate \eqref{art:matched-local-model} with
$\lambda=a_\ell^2$, $\sigma=a_\ell^4$ and $C=a_\ell^3$. The containment just
proved applies before each candidate is tested, so the scaling argument
transfers every bracket decision. Equivalently, the additive constant
cancels in Armijo and every function difference and gradient--step product
is the seed quantity times $a_\ell^4\rho^m>0$. The actual search traverses
exactly the candidate list and makes the same strict decisions.
In \eqref{art:failure_update}, the projection factors are unchanged and
the rank-one term has factor $\rho^m$. The cycle closures prove the next
point and matrix identities, completing the induction.

The finite matrices are positive definite. Equation~\eqref{art:upper_orbit}
gives convergence to $\bar x_\ell$ and strict descent to $a_\ell^3>0$.
Since the base query set is finite, the norm of every query is bounded by
\[
 a_\ell\|(2,-32)\|+a_\ell^2\max_{\text{base points and queries}}\|z\|,
\]
which tends to zero. The starting points also tend to zero. All other
claims were proved directly from \eqref{art:upper_objective}.
\end{proof}


Figure~\ref{art:shrinking-scales} makes the local quantifier explicit:
one function contains every shrinking copy. For each fixed copy the run
converges to its own wrong center, while the starting points and all their
queries can be made arbitrarily close to the same true minimizer.

\begin{figure}[t]
\centering
\includegraphics[width=\textwidth]{figures/shrinking_scales.pdf}
\caption{The $C^{2,1}$ local embedding, schematically. In one fixed
objective, bad centers $\bar x_\ell\to x_*$ have affine-model
neighborhoods of width $O(a_\ell)$. Each complete copied query orbit fits
inside a region of size $O(a_\ell^2)$ and converges to $\bar x_\ell$.
The inset shows the seed and two contraction copies as a finite illustration
of the infinite induction. Locations, neighborhood shapes and relative
scales are schematic, not a plot of the infinitely defined objective.
The lower-curvature construction uses the different spatial factor
$a_\ell$, as explained below.}
\label{art:shrinking-scales}
\end{figure}

\subsection{What changes when positive lower curvature is required}
The flat affine regions in the first construction explain why its convex
branches have no uniform positive lower curvature. Adding the same positive
quadratic to each branch restores that curvature, but changes the dynamics.
The old exact cycle can no longer be copied without adjustment. The next result constructs
a nearby trajectory by a stable-manifold argument.
\par\medskip
\noindent\begin{minipage}{\textwidth}
\begin{theorem}[Independent local failure with positive lower curvature]
\label{art:local-failure-lower}
There is a finite strongly convex function $F:\R^2\to\R$ with compact
sublevels and a unique nonsmooth minimizer at zero such that the following
holds. For the fixed deterministic unit-first bracketing implementation of
Section~\ref{art:failure}, with $c_1=10^{-4}$, $c_2=1/2$ and $H_0=I$,
there are starting points $x_0^{(\ell)}\to0$ whose infinite full-BFGS runs
converge to nonstationary points $\bar x_\ell\ne0$, with $\bar x_\ell\to0$ and
$F(x_k^{(\ell)})\downarrow F(\bar x_\ell)>F(0)$. Every accepted and rejected query
is differentiable, the accepted lengths repeat $1,1/4,1/16$, and the
supremum of all query norms in the $\ell$th run tends to zero.

The objective can be a maximum of four globally $C^{1,1}$, uniformly
strongly convex branches with globally Lipschitz gradients. There are
fixed $0<\underline\mu\le\overline L<\infty$ such that
\[
 \underline\mu I\preceq\nabla^2F(x)\preceq\overline L I
 \quad\text{at every twice-differentiable point}.
\]
The first three branches in this construction are not twice
differentiable at zero.
\end{theorem}

\end{minipage}
\par


\begin{table}[t]
\centering
\begin{tabular}{p{.29\textwidth}p{.27\textwidth}p{.32\textwidth}}
\hline
Branch regularity & Uniform curvature & At the accumulation point\\
\hline
Four convex $C^{2,1}$ branches
& $0\preceq\nabla^2f_i\preceq64I$ globally
& All branch Hessians exist and are Lipschitz; no uniform positive lower bound\\[3pt]
Four uniformly strongly convex $C^{1,1}$ branches
& Positive lower bound and globally Lipschitz branch gradients
& The first three branches are not twice differentiable\\
\hline
\end{tabular}
\caption{The two independent local failure regimes. Each uses one fixed
objective and $H_0=I$, with bad starts arbitrarily close to its unique
minimizer; all actual accepted and rejected queries are differentiable
and remain in neighborhoods shrinking to that minimizer. The respective
line-search parameter ranges are stated in
Theorems~\ref{art:local-failure} and~\ref{art:local-failure-lower}.}
\label{art:regularity-map}
\end{table}

Here a classical Hessian is the derivative of the gradient map, with that
map defined in a neighborhood of the point. A pointwise second-order
expansion alone is not the convention.

The proof in Appendix~\ref{art:lower_curvature} has the same sequence of
tasks as the first construction. It starts from the explicit strongly convex
objective $Q_\mu(z)=\mu\|z\|^2/2+G(z)$, whose true minimizer is away from
the wrong corner. A normalized three-step return map has the original seed
as a hyperbolic boundary fixed point. Its stable manifold contains positive
scale states, yielding an actual perturbed bad orbit. The seed's strict
finite margins ensure that the same accepted and rejected trial lists
persist along that orbit. No rounded initial state is asserted to lie on
the stable manifold.

For the local embedding the scalar node values are now $a_\ell^2$ and the
node slopes $a_\ell$. The resulting model is
$F(\bar x_\ell+a_\ell z)=a_\ell^2(C_*+Q_\mu(z))$, so
$\lambda=a_\ell$, $\sigma=a_\ell^2$ and again
$\lambda^2/\sigma=1$. We choose one entire perturbed orbit small enough
for every query to fit the affine scalar regions. Its one initial SPD
matrix is then used in every copy; one fixed affine change makes it $I$.
The faster variation of the scalar derivative retains global $C^{1,1}$
regularity but prevents second differentiability at the accumulation point.
This is why the theorem does not supply a finite strongly convex $C^2$
local counterexample.

The two constructions separate upper branch regularity from positive lower
curvature. Their common mechanism is a local copy of an actual bad query
tree; the different scalar scales determine which regularity survives at
the accumulation point. The remaining convergence questions are discussed
next.
